\section{Discussion}
\label{sec:discussion}

\subsection{Answers to the Research Questions}

\textit{RQ1: do the skills mitigate forgetting? Supported on CartPole, for the tested sequences.} In the ablation, the variants that learn Mars and Jupiter are those with Plasticity. Among them, only ENA-Combined maintains performance in both environments at test, and it is the most reliable variant (Section~\ref{sec:ablation}). Adding the archive also changes how the test policy is chosen, so the ablation does not isolate the archive's contribution from the selection's. Without Plasticity, ENA fails to acquire competent policies on Jupiter. The reliability gains observed in the ablation therefore reflect both policy acquisition and subsequent retention. The single-policy learners degrade markedly on Jupiter, the CartPole environment trained just before the last, though they retain performance on Mars. On HalfCheetah, they keep \Res{cheetah/SAC/retention/left/mean}--\Res{cheetah/PPO/retention/left/mean}\% on average of what they reached on the two environments left behind, whereas ENA retains full performance across all previously trained dynamics (Sections~\ref{sec:forgetting_cartpole} and~\ref{sec:cheetah}). All of these findings hold for the specific implementations, task sequences, and budgets evaluated; alternative settings were not tested.

\textit{RQ1a: is an archive enough by itself? No.} Variants with a Hall of Fame but no detection are not more reliable than the same variants without it, and ENA-HoF is less reliable than ENA-NoSkill. Variants whose archive is organized by detected physics are more reliable than those with an archive but no detection, in both cases. With Plasticity, they are also more reliable than the variant without any archive, and only this combination maintains performance across all learned environments (Section~\ref{sec:ablation}). In this ablation, merely archiving proven specialists proved insufficient: archiving increased reliability only when indexed by the physical regime under which each specialist was validated. The recall rule also matters. Bayesian policy reuse is more reliable than deploying the recognized niche's highest-Trust gladiator. Niche match, in turn, is not significantly more reliable than ignoring the niche (Section~\ref{sec:selection_ablation}). With either simpler rule, the archive is not more reliable than no archive at all.

\textit{RQ1b: does the memory cost extra interaction? Partly answered.} Without Plasticity, archiving and detecting add no measurable candidate evaluation of their own. Their only measured cost is wall-clock time (Section~\ref{sec:ablation}). With Plasticity, adding the Hall of Fame comes with more candidate steps. The increase is not significant, but the data cannot exclude it, up to \Res{ablation/ENA-Plasticity-vs-ENA-Plasticity+HoF/compute/candsteps/ci/high} steps per replicate. The archive takes no extra candidate evaluation to fill, but it may change how much the failure-driven search spends.

\textit{RQ1c: can the memory be built online, without labels? Yes in both domains, with different drivers.} On CartPole, the transition-model detector separates the training environments, with the fragmentation and false alarms of Table~\ref{tab:detection}. On HalfCheetah, the same kind of detector fails (Section~\ref{sec:chirp}). The chirp, whose response does not depend on the agent, meets the detection target on the benchmark. It also gives each training environment its own niche in every replicate of the HalfCheetah study (Section~\ref{sec:detection_benchmark}). Neither driver uses a gradient.

\textit{RQ2: does failure-driven search pay? Partly supported.} The comparisons that add Plasticity show the largest differences in the ablation, and only the variants with Plasticity learn the harder environments (Section~\ref{sec:ablation}). In the ablation, however, the budget is a number of episodes, so more search also means far more interaction with the environment. Searching at full effort in every episode costs an order of magnitude more candidate evaluation (Section~\ref{sec:compute}), a cost set by construction. Whether the ramp reaches what full effort does remains open. In the single replicate run at full effort, full effort was the more reliable. Whether scaling search effort dynamically with failure outperforms distributing identical effort uniformly across iterations remains an open empirical question. The quadratic controller was more reliable than the ramp and used fewer candidate steps, but neither difference is significant (Section~\ref{sec:controller}). These findings therefore support the general principle of scaling search with failure, rather than the specific linear ramp evaluated here.

\subsection{How Detection May Help the Archive}
The outcome of RQ1a is interpreted as a limitation of recall rather than of storage capacity. The selection ablation is consistent with this reading, but neither ablation tests it directly. Without niches, the archive is one pool of gladiators proven in different physics. During training, their protected copies hold population slots, whatever the current physics. At test time, the selection mechanism lacks dynamical evidence to distinguish specialists evolved for Jupiter from those evolved for Earth. Adding the archive also replaces a random population member by Bayesian policy reuse (Section~\ref{sec:variants}). With detection, every gladiator carries the physics it was proven in. The dynamics likelihoods of the selection~\eqref{eq:bpr} point it to the matching niche, and the gladiators of one niche can no longer be evicted by another niche's.

Selection rules that use the same niches also differ in reliability. Bayesian policy reuse differs from niche match in two ways. It weighs each episode's dynamics likelihoods into a belief, instead of following the detector's single decision. It also learns from the test returns. The selection ablation does not separate the two. One plausible mechanism, not directly isolated in the experiments, involves the feedback provided by episode returns. They let Bayesian policy reuse abandon a gladiator that fails in its own niche, which niche match cannot do, and across a test phase, niche match evaluates substantially fewer distinct individuals (Section~\ref{sec:selection_ablation}). A gladiator archived for solving a niche could still fail there at test, and without return-based evidence, an ineffective specialist would continue to be deployed.

\subsection{Retention, Adaptation and Absolute Return}
\label{sec:retention}
ENA reliably preserves acquired behaviors, but exhibits lower sample efficiency and asymptotic return than single-policy learners in other respects, by a margin that varies with the domain. Both gaps match known limits of evolutionary search. Evolutionary algorithms need many samples, and struggle on problems that require optimizing many parameters~\cite{khadka2018evolutionary,pourchot2019cemrl}. The cost of exploring in parameter space grows with the number of parameters~\cite{vemula2019contrasting}, and random mutations tend to disrupt large networks~\cite{lehman2018safe}. This is why ENA's specialists are small. On CartPole, ENA reaches criterion performance in a new environment later than PPO, as the majority of environment steps are spent evaluating candidate perturbations prior to identifying a successful policy (Section~\ref{sec:forgetting_protocol}). SAC-Discrete scores slightly higher on the environments it does not forget, whereas ENA preserves the earlier environment on which single-policy baselines catastrophically degrade.

On HalfCheetah, the gap is larger. ENA preserves performance from earlier phases, but achieves lower asymptotic returns than SAC within the allocated step budget. The robot has six joints to control. PPO and SAC control them with hidden layers of \Res{cfg/cheetah/baselinenet} units, while ENA's specialists have \Res{cfg/cheetahena/brain} units. Evolving networks as large as the baselines' would mean searching far more parameters, at the cost described above. Even with small networks, most of ENA's steps go to evaluating candidates rather than to its acting specialist's episodes (Section~\ref{sec:forgetting_protocol}). What SAC loses on the environments left behind exceeds the return ENA reached there at the end of each phase (Section~\ref{sec:cheetah}). Most of ENA's retention above 100\% comes from its end-of-phase reference, which includes episodes of specialists outside the archive. Measured against its gladiators' episodes alone, its retention is \Res{cheetah/archived/ENA-Combined/retention/left/mean}\% (Section~\ref{sec:cheetah}).

A specialist archive is hypothesized to provide the greatest advantage when environments recur and relearning is computationally expensive, though this setting was not explicitly evaluated. When an environment is unlikely to return, or a gradient learner can relearn it quickly, a single strong policy that is partly forgotten may be worth more than a weaker one kept intact.

\subsection{Scope of the Comparison}
\label{sec:scope}
The single-policy learners measure how much standard learners forget on these streams. The claim of this paper rests on the skill ablation. The comparison has three limits. First, at test, every agent's weights are frozen, but ENA still routes among its archived specialists with an updating belief. The single-policy learners deploy one policy. The comparison thus contrasts context-indexed specialist preservation with single-policy learning, not two equally equipped systems.

Second, ENA evaluates its candidates in fresh copies of the current environment, which the experiment builds from the environment's own parameters (Appendix~\ref{app:information}). Every candidate step counts in ENA's budget. Still, the single-policy learners play only in the stream itself, and never use such a copy.

Third, the continual-RL methods closest to ENA are not baselines here, for reasons that differ between them. OWL~\cite{kessler2022same} and TeLAPA~\cite{lillo2026beyond} are told where tasks begin and end. OWL trains the head of the task it is given, and infers the task only at test. TeLAPA builds a task's archive, and retrieves a starting point from its archives, at each task boundary it is given. Both build their memory during the stream, but from labeled tasks. On the non-stationary streams evaluated here, these methods would require either explicit task boundaries—which are unavailable to any agent in this study—or an environment detector to infer them. Providing such an inferential mechanism would effectively endow those baselines with ENA's environment-detection capability. ITE~\cite{cully2015robots} builds its archive offline, before deployment. It adapts once to each damage, starting when performance drops below a user-set threshold. It neither keeps a memory across a sequence of conditions, nor recognizes a condition that returns.

RL-CD~\cite{dasilva2006dealing}, MBCD~\cite{alegre2021minimum} and SWOKS~\cite{dick2024statistical} do detect contexts online without labels, and keep a policy per context, so they address the same problem. A comparison with them, however, would be confounded. RL-CD's models are tabular and need a discretized state space, while MBCD and SWOKS learn by gradient descent. A difference could therefore come from the learning algorithm rather than from the memory. Direct comparisons with such algorithms are deferred to future investigations of hybrid architectures combining ENA's modular skills with gradient-based optimization (Section~\ref{sec:conclusion}). One promising formulation would train specialist policies per niche via gradient descent while retaining ENA's environment detection and Bayesian policy reuse.

Finally, wall-clock execution times reflect specific software implementations, measured on a single workstation for CartPole and on dedicated cloud instances per replicate for HalfCheetah. While environment steps and candidate steps represent intrinsic properties of the algorithms, wall-clock time also reflects implementation and hardware utilization efficiency; hence, it does not imply that ENA is inherently less computationally demanding across all execution platforms.

\subsection{Limitations}
\begin{itemize}
\item \textit{Two benchmark families.} Both use related physics changes, with a single morphology in each. Visual, partially observable or structurally different tasks are untested.
\item \textit{Fixed task order.} Each domain uses one training sequence (Mars--Jupiter--Earth; Earth--Heavy-Rough--Mars-Slippery). Ordering effects are untested.
\item \textit{No recurrence during training.} Earlier environments return only at the frozen test. Repeated switches, short dwell times and recall during training are untested.
\item \textit{Test-time choice in the ablation.} Without a Hall of Fame, there are no gladiators to choose from. In that case, ENA-NoSkill and ENA-Plasticity play a random member of their final population in every test episode. The variants with an archive choose among gladiators by Bayesian policy reuse. Adding the Hall of Fame thus changes both what is kept and how the test policy is chosen, and the ablation does not separate the two.
\item \textit{Interaction in the ablation.} The ablation fixes the number of episodes, not of environment steps. A variant with Plasticity takes many times more steps than one without (Section~\ref{sec:ablation}). The ablation shows what failure-driven search achieves with that interaction. It does not show that allocating effort by failure beats spreading the same effort evenly. Full search was run in a single replicate.
\item \textit{Resets.} The chirp requires resetting the environment at episode boundaries, to play its \Res{cfg/cheetah/chirp/steps}-step excitation from nominal initial states. Candidate evaluation needs more than a reset: a resettable copy of the current physics. Both are natural in simulation, digital twins and automated test rigs. They constrain continuous, non-resettable physical systems.
\end{itemize}
